\documentclass[11pt]{article}
\usepackage{amsmath,amssymb,amsthm}
\usepackage{geometry}
\usepackage{hyperref}
\usepackage{microtype}
\usepackage{booktabs}
\geometry{margin=1.2in}
\newtheorem{theorem}{Theorem}[section]
\newtheorem{corollary}[theorem]{Corollary}
\newtheorem{definition}[theorem]{Definition}
\newtheorem{remark}[theorem]{Remark}
\DeclareMathOperator{\logit}{logit}
\newcommand{\updateBelief}{\mathrm{updateBelief}}
\newcommand{\R}{\mathbb{R}}
\title{Applications of the Bayesian Network Interpretation of Transformers}
\author{Greg Coppola}
\date{2026}
\begin{document}
\maketitle
\begin{abstract}
The companion paper~\cite{coppola2026transformers} proves that a sigmoid
transformer can exactly realize belief propagation under a constructive
weight mapping. This paper develops the applications and implications of
that result. We organize them across six domains: theoretical foundations
(why LLMs work, scaling laws, formal reasoning), grounding and hallucination,
agents and planning, mechanistic interpretability, training objectives, and
efficiency. Each application is marked by its epistemic status --- formal
consequence, mechanistic interpretation, empirical hypothesis, architectural
proposal, or speculative extension --- so that the reader can calibrate
confidence accordingly. The paper concludes with an explicit account of
what the theory does not prove and a research roadmap for the open problems.
\end{abstract}
\tableofcontents
\newpage
\part{Introduction and Framework}
\section{Introduction}
\label{sec:introduction}

The companion paper~\cite{coppola2026transformers} proves that a sigmoid
transformer can exactly realize belief propagation (BP) under a constructive
weight mapping. Within the constructive setting, the forward pass implements
the BP update equations directly, with residual streams as belief states,
attention as evidence gathering, and feed-forward networks as belief updates.

The present paper asks: what follows from taking that result seriously? The
applications range from theoretical (a mechanistic account of why LLMs work)
to practical (grounding strategies, interpretability predictions, agent design)
to open research directions (calibrated uncertainty, factor graph recovery,
grounded planning systems).

\subsection*{Levels of Analysis}

The applications in this paper operate at several distinct levels that
should not be conflated:

\begin{center}
\begin{tabular}{lll}
\toprule
Level & Description & Example \\
\midrule
Exact constructive & Formally verified theorem & Sigmoid transformer realizes exact BP \\
Direct consequence & Follows from the theorem & Grounded BP permits verifiable correctness \\
Mechanistic interpretation & Component-level account & Attention = evidence gathering \\
Empirical hypothesis & Testable prediction & AND heads have rank-1 Q/K circuits \\
Architectural proposal & Design suggestion & Grounded transformer + QBBN hybrid \\
Speculative extension & Interpretive extrapolation & Planning as MDP inference \\
\bottomrule
\end{tabular}
\end{center}

Each section marks the epistemic status of its claims explicitly.

\subsection*{Map of Literatures}

Several independent research programs have converged on related pictures of
transformer computation. The BP result provides a common language for
organizing their findings:

\begin{center}
\begin{tabular}{lll}
\toprule
Literature & Main Idea & BP Connection \\
\midrule
Mechanistic interpretability & Circuits and routing heads & Message passing structure \\
In-context learning & Bayesian updating from examples & Posterior inference per forward pass \\
Graph neural networks & Node message passing & Explicit graph BP \\
Sparse autoencoders & Monosemantic semantic features & Latent factor graph nodes \\
RAG / tool use & External evidence injection & Conditioning on observed evidence \\
\bottomrule
\end{tabular}
\end{center}

\subsection*{What This Paper Is Not Claiming}

The formal theorem is about sigmoid transformers under a specific
constructive weight mapping. Production transformers use SwiGLU or GeLU,
not sigmoid. The theorem establishes that the architecture \emph{can}
implement exact BP; whether any particular trained model does is an
empirical question. The applications that follow are interpretations,
proposals, and hypotheses motivated by the result --- not derived
consequences of it, except where explicitly noted.
\part{Theoretical Consequences}
\section{Why LLMs Work: A Theoretical Foundation}
\label{sec:why-llms-work}

There is currently no widely accepted mechanistic theory of why next-token
prediction on text produces general intelligence. The standard story is
empirical: scaling laws hold, capabilities emerge, benchmarks improve.
Other theoretical programs exist --- information-theoretic accounts,
compression views, in-context learning theory, grokking --- but none has
achieved consensus as a mechanistic explanation of what the model is
computing.

\subsection*{The BP Interpretation}

If transformers implement belief propagation over an implicit factor graph,
then next-token prediction can be interpreted as marginal inference in a
Bayesian network over language concepts. Under this interpretation, the
model behaves as though it is learning a structured probabilistic model of
the world and performing inference in it.

This suggests a mechanistic account of what the model is computing: each
forward pass propagates beliefs through an implicit factor graph, updating
each proposition's probability given the evidence in the context window.
Next-token prediction is the output of this inference process --- the most
probable next token given the current belief state.

\subsection*{Scaling Through the BP Lens}

The scaling dimensions have natural BP interpretations:

\begin{center}
\begin{tabular}{ll}
\toprule
Dimension & BP Interpretation \\
\midrule
$W$ (sequence length) & Nodes in the factor graph \\
$L$ (layers) & Rounds of belief propagation \\
$H$ (attention heads) & AND patterns per round \\
$D_{ff}$ (FFN hidden units) & OR patterns per round \\
$D_{\mathrm{model}}$ (embedding) & Belief state capacity \\
\bottomrule
\end{tabular}
\end{center}

Scaling may improve performance because it expands inference capacity ---
more rounds, richer potentials, more simultaneous concepts. This is
qualitatively consistent with the BP interpretation, though the specific
power-law exponents are not derived from BP theory.

\subsection*{What This Does Not Explain}

The BP framework suggests a structure for what the model computes. It does
not explain why the training distribution (internet text) produces a useful
factor graph, nor does it derive the power-law scaling exponents from first
principles. These remain open questions.

\subsection*{Epistemic Status}

\begin{center}
\begin{tabular}{ll}
\toprule
Claim & Status \\
\midrule
Sigmoid transformers can implement BP & Formal theorem \\
Trained LLMs implement BP & Empirical hypothesis \\
Next-token prediction fits a Bayes net & Mechanistic interpretation \\
Scaling expands BP inference capacity & Mechanistic interpretation \\
BP explains power-law scaling exponents & Not established \\
\bottomrule
\end{tabular}
\end{center}
\section{Scaling Laws Have a Mechanistic Interpretation}
\label{sec:scaling-laws}

Kaplan et al.\ and subsequent work established that loss scales predictably
as a power law with compute, parameters, and data~\cite{kaplan2020scaling}.
This is an empirical regularity. No principled derivation of the exponents
has been established, though several theoretical programs have attempted
partial accounts.

\subsection*{The BP Interpretation}

Under the BP interpretation, scaling adds inference capacity in well-defined
ways: more layers add rounds of BP, more $D_{ff}$ adds OR gate capacity,
more $D_{\mathrm{model}}$ adds superposition capacity, and more $W$ adds
factor graph nodes. Each of these is a smooth, monotone increase in a
well-defined quantity. Power-law scaling in loss is qualitatively consistent
with a system where each additional unit of capacity provides diminishing
but consistent marginal improvement in inference quality over a fixed
distribution.

\subsection*{The AND Head Invariant}

Empirically, boolean-AND head count stays roughly fixed within a model
family while hub and OR capacity scales. This is consistent with the BP
interpretation: the routing primitive is compact and does not need to scale,
while inference capacity does. This evidence is preliminary --- three model
families --- and should not be overstated.

\subsection*{Epistemic Status}

\begin{center}
\begin{tabular}{ll}
\toprule
Claim & Status \\
\midrule
Scaling dimensions have BP interpretations & Mechanistic interpretation \\
More layers enable deeper reasoning chains & Mechanistic interpretation \\
AND head invariance is consistent with BP & Empirical observation (preliminary) \\
BP explains power-law exponents & Not established \\
\bottomrule
\end{tabular}
\end{center}
\section{Formal Reasoning as a Special Case}
\label{sec:formal-reasoning}

Classical first-order logic and probabilistic graphical models are usually
presented as separate formalisms. The QBBN~\cite{coppola2024qbbn} unifies
these: it is a probabilistic graphical model over propositions with AND/OR
boolean structure, proven consistent and complete with respect to the
first-order calculus. Logical deduction is the special case where the OR
gate weights are deterministic; statistical inference is the general case
where the weights are soft and learned from data.

\subsection*{The Controlled/Learned Spectrum}

This suggests a spectrum between controlled and learned reasoning:

\begin{itemize}
\item \textbf{Fully controlled}: OR gate weights are fixed to be
deterministic. The system implements exact logical deduction. This is the
QBBN in symbolic mode.

\item \textbf{Fully learned}: OR gate weights are learned from data. The
system implements soft probabilistic inference. This is the standard LLM.

\item \textbf{Hybrid}: some weights are fixed (known logical rules) and
some are learned (uncertain statistical relationships). This is the target
architecture for systems combining flexibility with formal reliability.
\end{itemize}

The transformer's BP structure is compatible with all three modes. Current
training produces the fully learned mode.

\subsection*{The Completeness Connection --- and Its Limits}

The QBBN paper proves consistency and completeness with respect to the
first-order calculus. This is a result about \emph{representational
expressivity}, not about practical inference capability. It says the QBBN
can in principle express any first-order fact --- not that a transformer
will reliably compute it, nor that inference will be tractable, nor that
the right weights will be learned from data.

Concretely: a transformer with grounded BP weights is not thereby guaranteed
to do mathematics reliably. Mathematical reasoning requires not just the
right representational structure but also the right weights, sufficient
layers for the required inference depth, and a training procedure that
produces calibrated beliefs. The completeness result says the architecture
is not the bottleneck. It does not say the other bottlenecks are easy.

\subsection*{Fast vs.\ Slow Reasoning}

The QBBN analysis connects to Kahneman's fast/slow
distinction~\cite{kahneman2011thinking}. Fast reasoning is forward
inference: causes propagate to effects in one pass of BP. Slow reasoning
requires reasoning by cases --- exploring multiple hypotheses, maintaining
multiple belief states, backtracking. The number of transformer layers
bounds the depth of fast reasoning; slow reasoning requires iteration or
explicit search outside the model.

\subsection*{Epistemic Status}

\begin{center}
\begin{tabular}{ll}
\toprule
Claim & Status \\
\midrule
QBBN is consistent and complete w.r.t.\ first-order logic & Formal theorem \\
Transformers are compatible with formal reasoning & Direct consequence \\
Trained transformers reliably do formal reasoning & Not established \\
Completeness implies practical mathematical competence & Explicitly not claimed \\
Hybrid architecture is achievable & Architectural proposal \\
\bottomrule
\end{tabular}
\end{center}
\part{Grounding and Hallucination}
\section{Hallucination Has a Formal Definition}
\label{sec:hallucination}

Hallucination is typically described informally: the model says something
false or unsupported. This is a behavioral description, not a mechanistic
one. It gives no account of why hallucination happens or what would prevent
it.

\subsection*{The BP Definition}

The BP framework provides a formal definition: a model hallucinates when
its output beliefs diverge from the correct posterior given the evidence.
For a grounded transformer with a declared ontology and verifier, this
divergence is measurable:
\[
b(j) \neq P_{\mathrm{true}}(j),
\]
where $P_{\mathrm{true}}$ is the correct posterior induced by the knowledge
base and observed evidence.

For an ungrounded transformer, the comparison is not well-defined because
there is no formally declared ground truth. Hallucination in ungrounded
models is not a numerical error in the inference procedure --- it is a
consequence of operating without a fully specified and verifiable grounding
structure.

\subsection*{What Grounding Enables --- and Its Limits}

A fully grounded transformer would need: a declared concept space, a
declared relational structure, a grounding procedure mapping natural
language to the declared system, and a verification mechanism.

Grounding does not eliminate all possible error. Even a formally grounded
system can have an incomplete ontology, incorrect premises, ambiguous
mappings from natural language, or approximation errors in inference.
What grounding provides is not a guarantee of truth but a guarantee of
\emph{verifiable correctness relative to a declared ontology}: outputs can
be checked against the knowledge base, and failures can be traced to
specific premises or inference steps. This is a weaker but more honest
claim --- and still substantially stronger than what ungrounded systems
can provide.

Retrieval-augmented generation is a partial form of grounding: it
temporarily constrains the inference space but does not provide persistent
ontological structure or a verifier.

\subsection*{Epistemic Status}

\begin{center}
\begin{tabular}{ll}
\toprule
Claim & Status \\
\midrule
Grounded BP permits formally verifiable correctness & Direct consequence \\
Grounding eliminates all hallucination & Explicitly not claimed \\
Ungrounded models lack a formal correctness criterion & Direct consequence \\
Hallucination requires grounding to formally address & Mechanistic interpretation \\
Scaling alone cannot eliminate hallucination & Mechanistic interpretation \\
\bottomrule
\end{tabular}
\end{center}
\section{The Grounding Gap}
\label{sec:grounding-gap}

There are two distinct kinds of reliability a language model can have.
\emph{Empirical reliability} means the model produces correct outputs most
of the time on a given distribution, as measured by benchmarks or human
evaluation. \emph{Formal reliability} means the model's outputs are
verifiable as correct relative to a specified knowledge base and ontology.

Current LLMs have empirical reliability. Grounded BP systems --- transformers
operating over a declared factor graph with a verifier --- have formal
reliability in restricted settings: outputs can be checked against the
declared knowledge base and failures traced to specific premises. The gap
between these two is the grounding gap.

\subsection*{Why Scaling Does Not Close It}

Scaling improves empirical performance on many benchmark distributions, but
the relationship between scale and reliability is task-dependent ---
calibration can worsen and failure modes can shift. More fundamentally,
scaling does not create a declared ontology, a finite concept space, or a
verifier. A 405B parameter model with no grounding structure cannot provide
formal correctness guarantees, regardless of its benchmark performance.

The grounding gap is not a quantitative gap that more parameters will close.
It is a structural gap between two different kinds of system.

\subsection*{What Would Close It}

Closing the grounding gap requires building systems that combine neural
inference (the transformer's BP computation) with symbolic grounding (a
declared ontology and verifier). This is the QBBN program. Whether it can
be scaled to general-purpose language understanding is the central open
problem.

\subsection*{Epistemic Status}

\begin{center}
\begin{tabular}{ll}
\toprule
Claim & Status \\
\midrule
Two kinds of reliability are distinct & Conceptual distinction \\
Grounded BP systems have formal verifiability & Direct consequence \\
Scaling reliably improves all reliability metrics & Explicitly not claimed \\
Scaling does not close the grounding gap & Mechanistic interpretation \\
QBBN hybrid closes the gap & Architectural proposal \\
\bottomrule
\end{tabular}
\end{center}
\part{Agents and Planning}
\section{Agents and the BP Vocabulary}
\label{sec:agents}

An LLM agent is typically described as: perception $\to$ reasoning $\to$
action, iterated. This description is functional but not formal. The BP
framework provides a formal vocabulary for describing each component ---
not a proof that agents are Bayesian networks in every mechanistic detail,
but a claim that the BP language maps naturally onto agent components and
may provide useful formal tools.

Under this interpretation: perception can be described as evidence
injection; reasoning as BP inference; action as decision inference over
the posterior; memory as persistent factor graph state; and tool calls as
factor graph queries. Whether these correspondences are exact or approximate
depends on how grounded the underlying factor graph is. For an ungrounded
LLM, they are interpretive. For a grounded transformer operating over a
declared QBBN, they become formal.

\subsection*{The QBBN as the Logic Layer}

The structured reasoning connecting perception to action can be formalized
as a QBBN with AND/OR factor structure. The LLM provides neural inference
capacity; the QBBN provides logical structure that makes reasoning
compositional and verifiable. This suggests a hybrid architecture where
every component is described in the same probabilistic language. Whether
this can be built and scaled is an open research question.

\subsection*{Epistemic Status}

\begin{center}
\begin{tabular}{ll}
\toprule
Claim & Status \\
\midrule
BP vocabulary maps onto agent components & Mechanistic interpretation \\
LLM reasoning is BP inference & Empirical hypothesis \\
Grounded agent loop has formal BP account & Architectural proposal \\
Formal verification of agent steps is possible & Speculative extension \\
\bottomrule
\end{tabular}
\end{center}
\section{Planning as Inference}
\label{sec:planning}

Planning is the problem of selecting actions to achieve goals. The standard
formal model is the Markov Decision Process (MDP): states $S$, actions $A$,
transition probabilities $P(s' | s, a)$, reward $R(s,a)$, and discount
$\gamma$. MDPs can be represented as dynamic Bayesian networks. Policy
evaluation is inference; policy improvement is optimization over the
inferred value function.

\subsection*{The Bellman Equation as Message Passing}

The Bellman equation
\[
V(s) = \max_a \left[ R(s,a) + \gamma \sum_{s'} P(s'|s,a)\, V(s') \right]
\]
can be interpreted as a message passing update: the value at state $s$ is
computed from the values of its neighbors, weighted by transition
probabilities. Under this interpretation, a transformer running BP over
an MDP factor graph computes value estimates iteratively across layers,
with each layer corresponding to one step of policy evaluation. This is
a mechanistic interpretation, not a proven equivalence for arbitrary
trained transformers.

\subsection*{Chain-of-Thought as Extended Inference}

When an LLM reasons about future states in a chain-of-thought, it is
plausibly approximating this process. One interpretation of why
chain-of-thought helps planning: it provides more effective rounds of
inference by externalizing intermediate estimates into the context window.
This is consistent with the BP interpretation but not uniquely predicted
by it.

\subsection*{Grounded Planning}

A grounded planning system would combine a formally declared state space
(QBBN propositions), learned or specified transition probabilities (OR gate
weights), a BP inference engine (the transformer), and a verifier. Whether
this architecture can be built and scaled is an open research question.

\subsection*{Epistemic Status}

\begin{center}
\begin{tabular}{ll}
\toprule
Claim & Status \\
\midrule
MDPs can be represented as dynamic Bayes nets & Established (standard formalism) \\
Bellman equation resembles BP message passing & Mechanistic interpretation \\
Chain-of-thought provides more inference rounds & Empirical hypothesis \\
Ungrounded LLMs fail at long-horizon planning & Empirical observation \\
Grounded transformer enables verified planning & Architectural proposal \\
\bottomrule
\end{tabular}
\end{center}
\section{Multi-Agent Systems and Distributed Inference}
\label{sec:multi-agent}

If each agent is performing BP over its local factor graph, then a
multi-agent system can be interpreted as a distributed inference computation
over a larger global factor graph. This resembles loopy belief propagation
on a graph too large to fit in a single model.

This is an interpretive framework, not a proven equivalence. Real LLM
agents are not calibrated BP nodes. Communication between agents is lossy
text, not exact belief messages. The convergence guarantees from formal
loopy BP do not transfer directly to multi-agent LLM systems.

\subsection*{What the BP Lens Suggests}

Even without exact equivalence, the distributed BP framing suggests useful
design principles: agent boundaries may work better when aligned with factor
graph structure; message formats carrying structured uncertainty may
coordinate better than raw text; convergence criteria inspired by BP belief
stability could inform stopping rules for multi-agent deliberation.

\subsection*{What Would Be Required for a Formal Account}

For the distributed BP interpretation to become a formal account, agents
would need to maintain calibrated beliefs over a shared declared factor
graph, pass approximately correct belief messages, update beliefs according
to BP rules, and operate over a graph structure where loopy BP converges.
None of these conditions hold for current LLM multi-agent systems.

\subsection*{Epistemic Status}

\begin{center}
\begin{tabular}{ll}
\toprule
Claim & Status \\
\midrule
Multi-agent systems resemble distributed BP & Mechanistic interpretation \\
BP lens suggests useful design principles & Architectural proposal \\
Current LLM agents converge to correct posteriors & Not supported \\
BP-inspired coordination protocols may help & Speculative hypothesis \\
\bottomrule
\end{tabular}
\end{center}
\part{Interpretability}
\section{Interpretability as Hypothesis Testing}
\label{sec:interp-hypothesis}

Mechanistic interpretability is currently largely exploratory. The BP
framework changes this: it generates specific, testable predictions about
what trained transformers should look like at the weight level.

\subsection*{Specific Predictions}

\begin{itemize}
\item Boolean-AND heads should have approximately rank-1 Q/K circuits
consistent with \texttt{projectDim}-style matching.
\item The dominant singular vector of the Q/K circuit should correspond to
a monosemantic SAE feature.
\item The V circuit should show \texttt{crossProject}-style routing from a
belief-encoding dimension to a scratch-slot dimension.
\item $W_1$ rows of the FFN should cluster by input SAE feature.
\item $W_2$ columns should correspond to interpretable belief updates.
\item The relationship between input and output SAE features for each FFN
unit should be interpretable as a conditional probability table entry.
\end{itemize}

These are specific structural claims that can be confirmed or falsified by
inspecting the weights of a trained model with a good SAE.

\subsection*{What a Negative Result Would Mean}

If boolean-AND heads do not show rank-1 Q/K structure, that would falsify
the constructive BP account for those heads --- not just complicate it. The
BP framework is a genuine scientific hypothesis because it makes predictions
that could be wrong.

\subsection*{The Closing Experiment}

Take a model with a trained SAE. Identify boolean-AND heads by behavior.
Project the Q/K matrix onto the SAE feature basis. Check if the dominant
component corresponds to a monosemantic feature. Check if the V circuit
routes from a belief dimension to a scratch slot. Check if input/output SAE
feature pairs for active FFN units are interpretable as conditional
relationships. If this succeeds, the implicit factor graph is legible.

\subsection*{Epistemic Status}

\begin{center}
\begin{tabular}{ll}
\toprule
Claim & Status \\
\midrule
BP predicts rank-1 Q/K structure in AND heads & Empirical hypothesis \\
SAE features correspond to factor graph nodes & Empirical hypothesis \\
Factor graph is legible from weights + SAE & Empirical hypothesis \\
Negative SAE result would falsify BP account & Direct consequence \\
\bottomrule
\end{tabular}
\end{center}
\section{SAE Features as Factor Graph Nodes}
\label{sec:sae-features}

FFN hidden units are not the right unit for understanding the implicit
factor graph. A single hidden unit may participate in multiple factor
relationships via superposition, and a single factor relationship may
require multiple hidden units to implement accurately with non-sigmoid
activations.

The natural candidates for factor graph nodes are SAE features.
Each monosemantic SAE feature --- one that activates for a single coherent
concept --- is a natural candidate for one node in the implicit factor
graph, with its activation value encoding the belief at that node. This
is a hypothesis, not a proven correspondence: SAE features may themselves
be mixtures, basis artifacts, or local linear approximations rather than
clean factor graph nodes.

\subsection*{Toward Partial Factor Graph Recovery}

The BP interpretation suggests a pathway toward partial factor graph
recovery. Given a sufficiently complete SAE, factor graph nodes would
correspond to monosemantic SAE features; edges would be identifiable from
$W_1$ rows that respond to pairs of input features whose $W_2$ columns
write to output features; and factor potentials would be the learned
weights connecting input to output features.

Current SAEs cover a small fraction of the residual stream's
representational capacity. The factor graph may also be distributed,
redundant, or dynamically contextual in ways that make complete recovery
difficult even in principle. As SAE coverage improves, partial recovery
becomes more tractable. The BP result gives interpretability research a
specific target and a specific prediction to test.

\subsection*{Epistemic Status}

\begin{center}
\begin{tabular}{ll}
\toprule
Claim & Status \\
\midrule
SAE features are natural candidates for factor graph nodes & Empirical hypothesis \\
Monosemantic features correspond exactly to BP nodes & Not established \\
Partial factor graph recovery is possible in principle & Empirical hypothesis \\
Complete recovery is straightforward given better SAEs & Explicitly not claimed \\
Golden Gate Claude demonstrates evidence injection & Empirical observation \\
\bottomrule
\end{tabular}
\end{center}
\part{Training and Efficiency}
\section{Training Objectives Through the BP Lens}
\label{sec:training}

Next-token prediction optimizes cross-entropy loss on the training
distribution. This is equivalent to maximum likelihood estimation of the
implicit factor graph parameters. The bayes-learner experiment confirms
that gradient descent finds approximately correct BP weights from this
objective (val MAE 0.000752).

\subsection*{Alternative Objectives}

The BP framework suggests objectives that would produce better-calibrated
Bayes nets:

\begin{itemize}
\item \textbf{Calibration objectives}: penalize overconfident or
underconfident posteriors directly, not just cross-entropy.

\item \textbf{Grounding objectives}: reward outputs that are verifiable
against a declared knowledge base.

\item \textbf{BP consistency objectives}: penalize violations of the BP
fixed-point equations across layers. If beliefs at layer $l+1$ are
inconsistent with what BP would predict given beliefs at layer $l$,
penalize that divergence.
\end{itemize}

\subsection*{The Uniqueness Theorem as a Training Target}

The companion paper proves a uniqueness result: exact posteriors uniquely
force BP weights. This means there is a well-defined target for training
--- the weights that implement exact BP --- and the question is how well
current training objectives approach it.

\subsection*{Epistemic Status}

\begin{center}
\begin{tabular}{ll}
\toprule
Claim & Status \\
\midrule
Gradient descent finds approximate BP weights & Empirical observation \\
Calibration objectives would improve reliability & Architectural proposal \\
BP consistency training would enforce layer-wise calibration & Architectural proposal \\
Uniqueness theorem gives a well-defined training target & Direct consequence \\
\bottomrule
\end{tabular}
\end{center}
\section{Sparse Inference and Pruning}
\label{sec:efficiency}

The implicit factor graph is sparse: each node has only a few neighbors,
and most conditional relationships are zero or near-zero. If the sparse
factor graph structure can be identified --- which SAE features matter for
a given task, which AND connections are active, which OR gates are relevant
--- the rest can be skipped. This is structured pruning with a theoretical
justification.

\subsection*{Task-Specific Factor Graphs}

Different tasks activate different subgraphs of the implicit factor graph.
This is the theoretical foundation for mixture-of-experts and sparse
attention: they are approximations to task-specific factor graph routing,
and the BP framework gives a principled account of what they are
approximating.

\subsection*{Distillation as Factor Graph Compression}

Knowledge distillation can be understood as factor graph compression:
finding a smaller factor graph that approximates the posteriors of the
larger one. The BP framework gives a principled objective for distillation:
minimize KL divergence between the posteriors of the teacher and student
factor graphs, not just between their output distributions.

\subsection*{Epistemic Status}

\begin{center}
\begin{tabular}{ll}
\toprule
Claim & Status \\
\midrule
Implicit factor graph is sparse & Mechanistic interpretation \\
Task-specific subgraph activation motivates sparse computation & Mechanistic interpretation \\
MoE approximates task-specific factor graph routing & Mechanistic interpretation \\
Distillation as factor graph compression & Architectural proposal \\
\bottomrule
\end{tabular}
\end{center}
\part{Limitations and Open Problems}
\section{What This Does Not Prove}
\label{sec:limitations}

The result that sigmoid transformers can exactly realize belief propagation
is formal and verified. The applications that follow range from
well-established to speculative. This section states clearly what the
theory does not prove.

\subsection*{The Theorem Is About Sigmoid Transformers}

Production transformers use SwiGLU or GeLU, not sigmoid activation. The
exact BP equivalence holds for sigmoid transformers under the constructive
weight mapping. For production transformers, the architecture is compatible
with BP-like computation and in synthetic experiments gradient descent
finds approximately correct BP solutions, but exact equivalence is not
established for arbitrary trained models.

\subsection*{Attention Is Not Literally Logical Conjunction}

The AND/OR framing is mechanistically motivated but is an interpretation,
not a theorem about arbitrary trained attention heads. Across three studied
model families (GPT-2 small, GPT-2 medium, Qwen~2.5 0.5B), a substantial
fraction of heads fall into the boolean-AND, continuous-OR, and hub
categories. The remainder exhibit context-dependent behavior not yet
characterized by the framework. Reproducibility and generalization to other
model families remain to be established.

\subsection*{The Factor Graph Is Not Directly Observable}

The implicit factor graph is encoded in learned weights, not explicitly
declared. Current SAEs cover a small fraction of the residual stream's
representational capacity. Factor graph edges have not been systematically
recovered. The claim that SAE features are natural candidates for factor
graph nodes is an empirical hypothesis with supporting evidence, not a
proven theorem.

\subsection*{What Would Falsify the Core Interpretation}

The BP interpretation would be significantly weakened if: boolean-AND heads
do not exhibit rank-1 Q/K structure under SAE analysis; SAE features fail
to organize into stable dependency structure; FFN updates are not locally
interpretable as belief updates; deeper layers do not improve iterative
reasoning behavior; or gradient descent does not find BP-like solutions on
clean synthetic tasks.
\section{Research Roadmap}
\label{sec:roadmap}

The BP interpretation of transformers is an active empirical program, not
a finished doctrine. The open research directions divide into three tiers
by tractability.

\subsection*{Tier 1 --- Most Directly Testable}

\textbf{Recover factor graphs from SAEs.} For each boolean-AND head, project
the Q/K matrix onto the SAE feature basis and check if the dominant
component corresponds to a monosemantic feature. Check if the V circuit
routes from a belief dimension to a scratch slot.

\textbf{Test layer-wise belief convergence.} Train a tiny transformer on a
synthetic Bayes net with known ground-truth posteriors. Measure how closely
the residual stream at each layer matches the BP belief after that many
rounds of message passing.

\textbf{Measure calibration by layer depth.} Compare output calibration of
models with different depths on multi-step reasoning tasks. If deeper models
are better calibrated on longer chains, that supports the
layers-as-BP-rounds account.

\subsection*{Tier 2 --- Requires New Infrastructure}

Complete SAE coverage of a production model. Build and benchmark a grounded
transformer + QBBN hybrid. Develop a calibration benchmark for BP
consistency that measures layer-wise belief stability.

\subsection*{Tier 3 --- Longer-Term}

Learning grounded factor graphs from text via expectation maximization.
Formal account of in-context learning as BP. Multi-agent BP convergence
conditions.

\subsection*{What Would Most Advance the Program}

In order of expected impact: (1) the SAE factor graph recovery experiment,
which closes the loop between the formal theorem and trained models; (2)
layer-wise convergence measurement, which tests the layers-as-rounds claim
directly; and (3) the grounded hybrid system benchmark, which demonstrates
the practical value of grounding.
\bibliographystyle{plain}
\bibliography{refs}
\end{document}